\itemhead{习题2.2 · 第3题（教材p.49）}{XJ-B1-C02-HW-P049-Y0622}
\label{XJ-B1-C02-HW-P049-Y0622}
{\footnotesize 湘教版教材p.49 / PDF第56页 / 习题2.2 / 第3题}
\field{原题}{\subq{1}已知二次函数 $y=ax^2+bx+c$ 的零点为 2 和 3，求二次函数 $y=cx^2-bx+a$ 的零点。
\subq{2}已知二次函数 $y=x^2+2x+a$ 在区间 $[-2,1]$ 上的最小值大于 0，求该函数的零点个数。}
\field{参考答案}{（1）原式比例关系 $b=-5a,c=6a$，新函数为 $a(6x^2+5x+1)$，零点 $-1/2,-1/3$；（2）顶点在区间内，$a-1>0$，故全域上均为正，无零点。}
{\small \field{关联知识}{Q01 二次函数、零点、判别式与根系关系（教材帮p.65 / PDF74）}
\field{方法与易错点}{按已知零点、顶点或系数选择交点式、顶点式或一般式；根平方和用根系关系消元。 易错：遗漏 $a
e0$；把二重根算成两个不同零点；有两正根只查积正。}
\field{答案依据}{答案册相关题组：PDF第13、14页（册内页12、13）；按当前题库原题逐项重算，不按旧版题号机械对套。}}
